\bookchapter{บทที่ 59: รู้จัก iOS — ระบบนิเวศของ Apple กับระบบนิเวศนักพัฒนา}

\chapterintro{ยินดีด้วยที่ตัดสินใจแบบนี้ — จาก “คนที่ใช้ iPhone” มาเป็น “คนที่เขียนแอปให้ iPhone” ก่อนลงมือ เรามากางแผนที่ดูกันก่อน: iOS มาจากไหน? “สวนกำแพง” ของ Apple แปลว่าอะไรสำหรับนักพัฒนา? แอปหนึ่งตัว จากนิ้วที่พิมพ์คีย์บอร์ดของคุณ ถึงเครื่องของผู้ใช้ ระหว่างทางเกิดอะไรขึ้นบ้าง?}

\booksection{ชาติกำเนิดของ iOS: โทรศัพท์เครื่องเดียวพลิกทั้งอุตสาหกรรม}

มกราคม 2007 สตีฟ จอบส์ดึง iPhone รุ่นแรกออกมาโชว์บนเวทีแถลงข่าว ตอนนั้นมันยังไม่รองรับแอปจากผู้พัฒนาภายนอกด้วยซ้ำ — จนกระทั่งปี 2008 Apple เปิดตัว App Store และ iPhone SDK เปิดทางให้นักพัฒนาทั่วโลกเขียนซอฟต์แวร์ให้ iPhone วินาทีนั้น “เศรษฐกิจแอป” ก็ถือกำเนิด: ถึงวันนี้ App Store มีแอปหลายล้านตัว ส่วนแบ่งรายได้ที่จ่ายให้นักพัฒนาสะสมรวมเกินหลายแสนล้านดอลลาร์

ตัวระบบเองก็วิวัฒนาการ: แรกเริ่มชื่อ iPhone OS เปลี่ยนเป็น iOS ในปี 2010; เลขเวอร์ชันไต่ขึ้นถึง iOS 17 (2023) แล้วในปี 2025 Apple เปลี่ยนมาตั้งชื่อตามปี — iOS 26 เปิดตัวปี 2025 ส่วนเวอร์ชันปัจจุบันเดือนสิงหาคม 2026 คือ iOS 27 ชื่อจะเปลี่ยน แต่มีอย่างหนึ่งไม่เคยเปลี่ยน: นิสัยจ่ายเงินของผู้ใช้ iOS กับชื่อเสียงเรื่องคุณภาพซอฟต์แวร์ ถือเป็นดีที่สุดในระบบนิเวศมือถือที่ทุกคนยอมรับ — แอปเดียวกัน รายได้ฝั่ง iOS มักสูงกว่าฝั่ง Android อย่างชัดเจน

\begin{center}
\begin{tabularx}{\linewidth}{|X|X|X|}
\hline
\textbf{หมุดหมาย} & \textbf{ปี} & \textbf{เรื่องสำคัญสำหรับนักพัฒนา} \tabularnewline
\hline
App Store เปิดตัว & 2008 & เปิดระบบนิเวศพัฒนาบุคคลที่สาม \tabularnewline
Swift เปิดตัว & 2014 & ภาษาพัฒนา iOS ยุคใหม่ถือกำเนิด \tabularnewline
SwiftUI เปิดตัว & 2019 & เฟรมเวิร์ก UI แบบ declarative และเส้นหลักของหนังสือเล่มนี้ \tabularnewline
iOS 26 (ปีแรกของการตั้งชื่อตามปี) & 2025 & ตั้งชื่อตามปี; Apple Intelligence ถูกผนวกอย่างเต็มตัว \tabularnewline
iOS 27 & 2026 & เวอร์ชันใหม่ล่าสุดปัจจุบัน (ฐานอ้างอิงของหนังสือเล่มนี้) \tabularnewline
\hline
\end{tabularx}
\end{center}

\begin{center}\small\color{gray}บันทึกเหตุการณ์สำคัญของการพัฒนา iOS (ข้อมูลถึงสิงหาคม 2026)\end{center}

\booksection{สวนกำแพง: “กฎ” กับ “กำไร” ของระบบนิเวศ Apple}

ระบบนิเวศของ Apple มักถูกเรียกว่า “สวนกำแพง” (Walled Garden) — กฎเยอะ คุมเข้ม: แอปต้องผ่านการตรวจก่อนขึ้นสโตร์ พัฒนาต้องใช้ Mac กับ Xcode การเผยแพร่ก็เดินผ่าน App Store เป็นหลัก ฟังดูอึดอัด? แต่สำหรับนักพัฒนา กฎชุดนี้แลกมาด้วยผลตอบแทนเป็นเงินจริงก้อนโต:

\begin{itemize}
  \item อุปกรณ์กระจัดกระจายน้อยมาก: รุ่น iPhone มีแค่ไม่กี่รุ่น ขนาดหน้าจอกับเวอร์ชันระบบเป็นเอกภาพสูง — ปรับให้รองรับ 5 รุ่น ก็ครอบคลุมผู้ใช้ถึง 95\% (ฝั่ง Android อาจต้องปรับหลายสิบรุ่น);
  \item อัตราการอัปเดตเวอร์ชันระบบสูงมาก: ระบบใหม่เปิดตัวไม่กี่สัปดาห์ ผู้ใช้ก็อัปเดตกันถึงเจ็ดแปดสิบเปอร์เซ็นต์ API ใหม่ใช้ได้อย่างมั่นใจ ไม่ต้องไปดูแล “ระบบเก่าเมื่อห้าปีก่อน”;
  \item ผู้ใช้จ่ายเงินเก่ง: ผู้ใช้ App Store ยอมจ่ายค่าแอปกับการซื้อในแอปมากกว่า;
  \item การตรวจแม้จะเข้ม แต่สโตร์สะอาด: ของเลียนแบม ซอฟต์แวร์กวนโอนมีน้อย ผลิตภัณฑ์ดีจึงถูกมองเห็นง่ายขึ้น.
\end{itemize}

\begin{tcbwarn}{ข้อควรรู้: เคลียร์เรื่องค่าใช้จ่ายกันก่อน}
การพัฒนา iOS มีข้อกีดขวางแบบแข็ง: ต้องมี Mac (Xcode มีเฉพาะบน macOS) การขึ้นสโตร์ต้องใช้บัญชีนักพัฒนา 99 ดอลลาร์/ปี (บัญชีฟรีสำหรับนักพัฒนารายบุคคล ดีบักบนเครื่องจริงในเครื่องตัวเองได้ แต่ขึ้นสโตร์ไม่ได้) คิดเรื่องเงินลงทุนส่วนนี้ให้ขาดก่อนออกเดินทาง — แต่ข่าวดีคือช่วงเรียนรู้ใช้บัญชีฟรีก็เอาอยู่สบาย จะเล่าในบทที่ 60

\end{tcbwarn}

\booksection{เขียนด้วยภาษาอะไร: Objective-C หรือ Swift?}

ประวัติการพัฒนา iOS ก็มี “พระเอก” สองคน Objective-C คือเสาหลักรุ่นบุกเบิก — ใช้ตั้งแต่ปี 2008 ถึง 2014; Swift คือภาษาใหม่ที่ Apple เปิดตัวปี 2014 แล้วพัฒนาต่อเนื่องมาถึง Swift 6 (เปิดตัวปี 2024 เพิ่มการตรวจสอบ concurrency อย่างเข้มงวด) ถึงปี 2026 โปรเจกต์ใหม่ใช้ Swift กันแทบทั้งหมด

คำตรง ๆ สำหรับคนที่เริ่มจากศูนย์: เรียน Swift ไปเลย Objective-C ไม่ต้องแม้แต่จะหยิบมาดู — เว้นแต่คุณเข้าทำงานแล้วต้องดูแล legacy code เมื่อสิบปีก่อน ข้อดีของ Swift: ไวยากรณ์สบายตา (ไม่มีภาระเซมิโคลอนกับไฟล์ header ยุคภาษา C), กลไก Optional กำจัด “crash จาก null pointer” ตั้งแต่ระดับภาษา, ต่อกับ SwiftUI ได้เนียน และยังมี Swift Playgrounds บน iPad ให้เรียนไปเล่นไปไปด้วย

สไตล์ Objective-C (สมบัติโบราณ)

\begin{lstlisting}
NSString *name = @"สมชาย"; [name stringByAppendingString:@"สวัสดี"]; // วงเล็บเหลี่ยม + พอยน์เตอร์ + เซมิโคลอน ครบสามอย่าง
\end{lstlisting}

สไตล์ Swift (สมัยใหม่ สบายตา)

\begin{lstlisting}
var name = "สมชาย" name += "สวัสดี" // เข้าใจง่ายเหมือนภาษาพูดธรรมชาติ
\end{lstlisting}

\booksection{แอปหนึ่งตัวถือกำเนิดขึ้นได้อย่างไร?}

เหมือนทุกแพลตฟอร์ม การถือกำเนิดของแอป iOS ก็มีไปป์ไลน์เป็นของตัวเอง แต่ฉบับของ Apple มี “ด่านตรวจ” เพิ่มมาอีกสองด่าน:

ไอเดีย

ออกแบบ (UI+ฟังก์ชัน)

เขียนโค้ด (Swift + SwiftUI)

คอมไพล์และเซ็น

ทดสอบ (โปรแกรมจำลอง+เครื่องจริง)

รีวิวโดย Apple

เผยแพร่บน App Store

\begin{center}\small\color{gray}ภาพที่ 1: กระบวนการทั้งหมดของแอป iOS จากไอเดียถึงขึ้นสโตร์ (สังเกตขั้นตอน “ตรวจสอบ” ที่เพิ่มเข้ามา)\end{center}

แนวคิดสำคัญในไปป์ไลน์:

\begin{itemize}
  \item ลายเซ็น (Signing): แอป iOS ทุกตัวต้องถูก “เซ็น” ก่อนจึงจะติดตั้งได้ — พิสูจน์ว่า “แพ็กเกจนี้มาจากคุณและไม่ถูกแก้ไข” ช่วงเรียนรู้ Xcode จะเซ็นอัตโนมัติด้วย Apple ID ของคุณ ตอนขึ้นสโตร์ค่อยใช้ใบรับรองนักพัฒนา รายละเอียดอยู่ในบทที่ 67
  \item ซิมูเลเตอร์ (Simulator): iPhone เสมือนที่รันบน Mac เป็นสมรภูมิหลักของการพัฒนา-ดีบัก มันใช้ CPU ของ Mac จึงลื่นพิเศษ ประสบการณ์ดีกว่าซิมูเลเตอร์ Android อยู่หลายขุม
  \item TestFlight: เครื่องมือ “แจกเวอร์ชันทดลองภายใน” ทางการของ Apple — ก่อนขึ้นสโตร์ ส่งเวอร์ชันทดสอบให้ผู้ใช้ทดสอบได้ถึง 10000 คน แล้วเขาติดตั้งลองใช้ผ่านแอป TestFlight นี่คือความเนี้ยบที่มีเฉพาะในระบบนิเวศ iOS
\end{itemize}

\booksection{ภาพรวมเทคโนโลยี iOS: “คลังอาวุธ” ของคุณ}

\begin{center}
\begin{tabularx}{\linewidth}{|X|X|X|}
\hline
\textbf{ชั้น} & \textbf{ตัวเลือกหลัก (2026)} & \textbf{ใช้ทำอะไร} \tabularnewline
\hline
ภาษา & Swift 6 & ภาษาสำหรับเขียนลอจิก \tabularnewline
หน้าจอ & SwiftUI + คู่มือออกแบบของ Apple (HIG) & วาดหน้าจอด้วยโค้ด (เส้นหลักของหนังสือ) \tabularnewline
สถาปัตยกรรม & MVVM + @Observable & แยกหน้าจอกับลอจิกออกจากกัน \tabularnewline
จัดเก็บข้อมูล & SwiftData (ฐานข้อมูล) + UserDefaults (ค่าตั้ง) & เก็บข้อมูลลงเครื่อง \tabularnewline
เครือข่าย & URLSession + async/await + Codable & คุยกับเซิร์ฟเวอร์ \tabularnewline
เครื่องมือบิลด์ & Xcode Build System / xcodebuild & แปลงโค้ดเป็นไฟล์ติดตั้ง \tabularnewline
สภาพแวดล้อมการพัฒนา & Xcode (มีเฉพาะบน macOS) & IDE ทางการของ Apple \tabularnewline
\hline
\end{tabularx}
\end{center}

\begin{center}\small\color{gray}ภาพรวมเร็วเทคโนโลยีหลักของ iOS ปี 2026\end{center}

สังเกตเห็นไหม: ลิสต์ชุดนี้สั้นกว่าฝั่ง Android อยู่หนึ่งตอน — เพราะ Apple จัดการเรื่องเยอะแยะให้ “ครบจบ” หมดแล้ว ราคาของระบบปิดคือความอิสระ สิ่งที่ได้กลับมาคือความสบายใจ: ไม่ต้องลังเลว่าจะใช้ไลบรารีเครือข่ายตัวไหน ฐานข้อมูลตัวไหน ของที่ทางการให้มามักเป็นคำตอบที่ดีที่สุดอยู่แล้ว

\begin{tcbpractice}{ลงมือทำ: ดู “ทะเบียนบ้านฝั่ง Apple” ของเครื่องคุณ}
เปิด [การตั้งค่า] [ทั่วไป] [เกี่ยวกับเครื่อง] ของ iPhone แล้วจดสามอย่าง: เวอร์ชัน iOS, ชื่อรุ่นเครื่อง, พื้นที่ว่างที่ใช้ได้ แล้วแวะเล่นหน้า “วันนี้” ใน App Store ตามอำเภอใจ สังเกตภาพหน้าจอกับคำโปรยของแอปคุณภาพดี — ตั้งแต่วันนี้ สายตาตอนเดินสโตร์ของคุณต้องเปลี่ยนจาก “ผู้ใช้” ไปเป็น “นักพัฒนา” แล้ว

\end{tcbpractice}

\begin{tcbsummary}{สรุปบทนี้}
\begin{itemize}
  \item ปี 2008 App Store เปิดตัว “เศรษฐกิจแอป” ถือกำเนิด; ปี 2026 เวอร์ชันปัจจุบันคือ iOS 27
  \item บัญชีของสวนกำแพง: กฎเยอะ แต่รุ่นอุปกรณ์น้อย เวอร์ชันใหม่ ผู้ใช้ยอมจ่าย สโตร์สะอาด
  \item ข้อกีดขวางแบบแข็ง: ต้องมี Mac; ขึ้นสโตร์ต้องใช้บัญชีนักพัฒนา 99 ดอลลาร์/ปี; ช่วงเรียนรู้บัญชีฟรีก็พอ
  \item ภาษาเรียน Swift ไปตรง ๆ; Objective-C เหลือไว้กับ legacy code เท่านั้น
  \item ไปป์ไลน์มีขั้นเพิ่มกว่า Android หนึ่ง: การตรวจของ Apple; TestFlight คือของวิเศษสำหรับแจกเวอร์ชันทดสอบ
  \item เทคโนโลยีมาตรฐานปี 2026: Swift 6 + SwiftUI + MVVM + SwiftData + URLSession + Xcode
\end{itemize}

\end{tcbsummary}
